\section{From living trees to commercial log supply}\label{from-living-trees-to-commercial-log-supply}

\subsection{1. Conditional McKean-Vlasov growth}\label{conditional-mckean-vlasov-growth}

Fix a biological mark class such as site, genetics, and management condition.
For a representative surviving tree, let

\[
\Xi_t=(s_{d,t},s_{h,t},q_t),
\]

where \texttt{s\_d=log\ d}, \texttt{s\_h=log\ h}, and \texttt{q\_t} contains continuous quality/form
states. Age follows

\[
da_t=dt.
\]

The paper uses the conditional McKean-Vlasov SDE

\[
d\Xi_t=b_G(\Xi_t,\hat\mu_t,a_t,\varsigma,a_t^G)dt
+\sigma_G(\cdot)dB_t
+\sigma_G^0(\cdot)dW_t^0,
\]

with conditional law

\[
\hat\mu_t=\mathcal L(\Xi_t\mid\mathcal F_t^0).
\]

Dynamic programming for controlled McKean--Vlasov systems with common noise is
developed by \citet{pham2017dynamic} and \citet{djete2019dynamic}. The forestry
model uses their conditional-law state formulation.

\texttt{W\^{}0} is common noise and \texttt{B} is idiosyncratic noise. Conditional on the common
noise, the population law remains random but the common environmental path is
shared.

This distinction matters in simulation. If a drought shock is intended to be a
site-year event, independently drawing a new drought shock for every tree would
artificially remove within-site correlation.

\subsection{2. What the law dependence is doing}\label{what-the-law-dependence-is-doing}

A McKean-Vlasov model allows the drift or diffusion of a focal tree to depend on
the population distribution. A generic competition term is

\[
\mathcal C(x,\mu)=\int K_{\rm comp}(x,x')\,\mu(dx').
\]

This formulation contains several special cases.

\subsubsection{Scalar density}\label{scalar-density}

If \texttt{K\_comp(x,x\textquotesingle{})=c} is constant, competition depends only on total mass or
population density. This can be useful but is biologically coarse.

\subsubsection{Size-weighted mean field}\label{size-weighted-mean-field}

If

\[
K_{\rm comp}(x,x')=c\,D(x')^2,
\]

then competition is proportional to a basal-area-like moment of the stand.

\subsubsection{Local spatial competition}\label{local-spatial-competition}

If tree location is available,

\[
K_{\rm comp}(x,x')=w(\|\ell-\ell'\|)g(D(x),D(x')),
\]

can represent distance-decaying and size-asymmetric competition.

Density-only competition is a special case. The law-dependent formulation also
supports size-weighted and local spatial kernels when data identify them.

\subsection{3. Why log-size coordinates are useful}\label{why-log-size-coordinates-are-useful}

Writing diameter and height in logarithmic coordinates gives positivity after
exponentiation and allows additive Gaussian perturbations without creating
negative physical sizes.

If \texttt{Y\_t=log\ D\_t} follows

\[
dY_t=\mu_Y(Y_t,\cdots)dt+\sigma_Y(Y_t,\cdots)dW_t,
\]

then It??'s formula gives the corresponding diameter dynamics

\[
dD_t=D_t\left(\mu_Y+\frac12\sigma_Y^2\right)dt
+D_t\sigma_YdW_t.
\]

Thus even simple additive log-space noise becomes multiplicative noise in
physical diameter.

This parametrization preserves positivity. Its empirical use requires evidence
for a Gaussian diffusion approximation.

\subsection{4. Sources of uncertainty}\label{sources-of-uncertainty}

The biological parametrization is left open. The market layer can be driven by a
deterministic increment model, a learned
transition density, a jump process, or a size-structured PDE.

When diffusion is used, separate at least three uncertainty sources:

\begin{enumerate}
\def\labelenumi{\arabic{enumi}.}
\tightlist
\item
  \textbf{process heterogeneity} - unobserved micro-site or biological variation;
\item
  \textbf{common environmental shock} - climate or disturbance shared across trees;
\item
  \textbf{measurement error} - error in observing DBH, height, or quality.
\end{enumerate}

Only the first two belong in the state process. Measurement error belongs in the
observation kernel.

Sparse remeasurement weakly identifies a complex state-dependent diffusion. The
examples therefore use simple noise for illustration.

\subsection{5. Well-posedness}\label{well-posedness}

Under standard Lipschitz dependence on state and conditional law in \texttt{W\_2},
linear growth bounds, and continuity in the control, the conditional
McKean-Vlasov SDE admits a unique strong solution on a finite horizon and its
conditional law remains in \texttt{P\_2}.

The coefficients satisfy a Lipschitz bound in state and population law:

\[
\|b(x,\mu)-b(x',\mu')\|
+\|\sigma(x,\mu)-\sigma(x',\mu')\|
\le L(\|x-x'\|+W_2(\mu,\mu')).
\]

This proposition establishes well-posedness for a specified growth law. Growth-law
identification requires longitudinal data.

\subsection{6. Finite-mass dynamics}\label{finite-mass-dynamics}

The normalized conditional law records composition and omits population mass.
Introduce harvest hazard \texttt{lambda\_t\^{}H(xi;u\_t)}, mortality hazard \texttt{lambda\_t\^{}M(xi)}, and a
recruitment/planting-rate measure \texttt{B\_t}.

For a smooth test function \texttt{phi}, the finite measure evolves weakly as

\[
d\langle\phi,\mu_t\rangle
=\left[\langle L_t\phi,\mu_t\rangle
-\langle(\lambda_t^H+\lambda_t^M)\phi,\mu_t\rangle
+\langle\phi,B_t\rangle\right]dt
+\text{common-noise term}.
\]

This is a transport-diffusion-reaction equation. Growth moves state; mortality
and harvest remove mass; recruitment adds mass.

\subsection{7. Harvest as a flow into log space}\label{harvest-as-a-flow-into-log-space}

Let \texttt{K\_B(dlambda\ \textbar{}\ xi,beta\_t\^{}B)} be a bucking kernel. Conditional on a harvested
Tree state \texttt{xi} and bucking rule \texttt{beta\_t\^{}B}, it returns a finite measure of logs.

Over a short decision interval, the expected log-count intensity is

\[
\Lambda_t(A)=\int_{E_T}
\lambda_t^H(\xi;u_t)K_B(A\mid\xi,\beta_t^B)\mu_t(d\xi).
\]

This formula is the core tree-to-log map.

It says that log supply depends jointly on:

\begin{itemize}
\tightlist
\item
  which trees exist;
\item
  which trees are harvested;
\item
  how each tree is cut;
\item
  the uncertainty encoded in the bucking/recovery kernel.
\end{itemize}

The tonne-weighted measure is

\[
\widetilde\Lambda_t(d\lambda)=m_L(\lambda)\Lambda_t(d\lambda).
\]

Again, no universal volume-to-mass coefficient is required.

\subsection{8. Bucking as a control problem}\label{bucking-as-a-control-problem}

Let \texttt{A\_B(xi)} be the feasible set of cutting patterns. Given a downstream
log-value function \texttt{v\_t\^{}L(lambda)} and bucking cost \texttt{c\_B}, define

\[
J_B(\xi,\beta;v_t^L)
=\int v_t^L(\lambda)K_B(d\lambda\mid\xi,\beta)-c_B(\xi,\beta).
\]

An optimal pattern satisfies

\[
\beta_t^{B,*}(\xi)\in\arg\max_{\beta\in A_B(\xi)}J_B(\xi,\beta;v_t^L).
\]

This creates a feedback from the market layer into physical recovery. If long,
high-grade logs are scarce and carry high downstream shadow value, the preferred
cut can differ from the cut that maximizes volume alone.

Under measurable-graph, compactness, measurability, and upper-semicontinuity
assumptions, a measurable maximizing selector exists.

\subsection{9. A simple competition example}\label{a-simple-competition-example}

Suppose log-DBH \texttt{Y\_i} follows

\[
dY_i=\left[r(1-e^{Y_i}/D_{\max})
-\gamma\,G(\mu_t)\right]dt+\sigma dB_i+\sigma_0dW^0,
\]

where

\[
G(\mu_t)=\int e^{2y}\,\mu_t(dy)
\]

is a basal-area-like moment.

This is a mean-field model: each tree responds to a competition pressure generated
by the whole size distribution. A stand with the
same stem density but many large competitors produces a larger \texttt{G(mu)}.

The stylized example shows how a size-weighted competition variable enters the
structural framework.

\subsection{10. Implementation}\label{implementation}

The market model takes a distribution of recoverable logs, with uncertainty, as
its input. Size classes, yield tables, and grade shares provide finite-dimensional
approximations to that distribution.
